\structelem{ВВЕДЕНИЕ}

Актуальность работы. Корпоративные сети Российской Федерации
остаются устойчивой целью несанкционированных воздействий,
включающих эксфильтрацию данных, латеральное перемещение,
разведку, отказ в обслуживании и эксплуатацию уязвимостей сетевых
сервисов. Доля шифрованного трафика в современных корпоративных
инфраструктурах превышает 90\,\%, что делает сигнатурный анализ
полезной нагрузки малоэффективным. Центр тяжести разработки систем
обнаружения вторжений (Intrusion Detection Systems, IDS)
переходит в сторону аномального анализа метаданных потоков с
применением методов машинного и глубокого обучения~\cite{khraisat2019survey,
ahmad2021nids}. При этом существенной проблемой остается
воспроизводимая, формализованная оценка таких систем в условиях,
приближенных к реальной эксплуатации, в том числе при дрейфе
распределений и вариативности атакующих сценариев.

Регулирование сетевой безопасности корпоративных инфраструктур
Российской Федерации опирается на Федеральный закон от 26.07.2017
№ 187-ФЗ <<О безопасности критической информационной инфраструктуры
Российской Федерации>>, Указ Президента Российской Федерации от
01.05.2022 № 250 <<О дополнительных мерах по обеспечению
информационной безопасности>>, Доктрину информационной
безопасности Российской Федерации (Указ Президента Российской
Федерации от 05.12.2016 № 646), приказ ФСТЭК России от 25.12.2017
№ 239 в действующей редакции, требующий внедрения средств
обнаружения и предотвращения компьютерных атак (группы требований
СОВ.1 и СОВ.2). Совокупность нормативных актов формирует объективный
спрос на отечественные методики и программные средства NIDS,
удовлетворяющие требованиям регулятора.

Объект исследования --- методы и технологии обнаружения аномальной
активности и несанкционированных действий в корпоративных сетях на
основе телеметрии сетевого трафика. Данные о поведении пользователей
рассматриваются в теоретической части как один из возможных
источников, в экспериментальной реализации и валидации не
используются.

Предмет исследования --- применение нейросетевых методов для
выявления аномалий и признаков несанкционированной активности по
данным сетевого трафика, а также разработка и применение
программного агента на основе конечного автомата и библиотеки
параметризованных атакующих шаблонов для синтеза трафика и
сценариев активного тестирования и валидации системы обнаружения.

Цель работы. Разработка методики обнаружения аномальной
активности и несанкционированных действий в корпоративной сети на
основе нейронных сетей, включающей анализ сетевого трафика,
классификацию угроз и моделирование атак с помощью программного
агента на основе конечного автомата, генерирующего трафик по
библиотеке параметризованных шаблонов и выступающего конечной
точкой воспроизводимого активного тестирования системы.

Задачи исследования:
\begin{enumerate}
    \item Аналитический обзор предметной области обнаружения
    аномалий и несанкционированных действий в корпоративных сетях;
    сравнительный анализ нейросетевых архитектур и публично
    доступных датасетов NIDS.
    \item Формализация функциональных и нефункциональных требований
    к системе обнаружения и определение количественных критериев
    качества.
    \item Проектирование методики и архитектуры нейросетевого
    детектора аномальной активности.
    \item Разработка подсистемы предобработки и формирования
    обучающих и тестовых выборок.
    \item Реализация и обучение нейросетевой модели детекции в
    едином программном регистре с совокупностью моделей-кандидатов.
    \item Разработка конечно-автоматного агента активного
    тестирования с библиотекой шаблонов и инжектором дрейфа.
    \item Реализация программного прототипа системы обнаружения
    в реальном времени с REST-интерфейсом и визуальной панелью.
    \item Экспериментальная проверка эффективности методики и
    подтверждение гипотезы исследования.
    \item Оценка эксплуатационных характеристик и оформление
    итоговых материалов.
\end{enumerate}

Гипотеза исследования. Применение нейронных сетей, обрабатывающих
последовательности агрегированных признаков сетевых потоков в
скользящих окнах, в сочетании с программным агентом на основе
конечного автомата, формирующего воспроизводимые сценарии
нормального и атакующего поведения по библиотеке параметризованных
шаблонов с контролируемым инжектированием дрейфа распределений,
обеспечивает построение методики обнаружения аномальной активности
и несанкционированных действий в корпоративной сети, превосходящей
классические сигнатурные и статистические подходы по комплексному
критерию качества, включающему точность и полноту классификации,
частоту ложных срабатываний при заданной полноте, задержку
детекции и устойчивость к контролируемому дрейфу распределений.

Методы исследования. Систематический анализ литературы;
математическое моделирование (вероятностные и
информационно-теоретические меры дрейфа, формальные определения
функций потерь и метрик); сравнительный анализ нейросетевых
архитектур и публично доступных датасетов NIDS по балльной шкале с
явно сформулированными критериями и весами; обучение нейронных
сетей и классических моделей в равных условиях; статистический
анализ результатов через bootstrap-доверительные интервалы и
парный критерий Уилкоксона.

Научная новизна. Новизна состоит в комбинации трех компонент:
(а)~интеграции нейросетевого детектора аномальной активности по
потоковому трафику, конечно-автоматного агента с библиотекой
параметризованных шаблонов и блока мониторинга дрейфа в единую
воспроизводимую методику; (б)~использования прозрачной
ground-truth по состояниям конечного автомата, обеспечивающей
измерение характеристик детектора, недоступных при использовании
только статических датасетов (per-state recall); (в)~единого
протокола экспериментов, включающего сравнение моделей в равных
условиях, анализ ошибок, калибровку порога под эксплуатационный
FPR, стресс-тестирование с агентом, drift-устойчивость и
эксплуатационную производительность.

Практическая значимость. Разработанный прототип реализован на
стеке Python\,/\,PyTorch с REST-интерфейсом на FastAPI и
визуальной панелью на Streamlit и может быть интегрирован в
подсистему сетевой аналитики корпоративного SOC. Зафиксированы
рекомендации по интеграции с источниками NetFlow\,/\,IPFIX,
калибровке порогов под допустимый FPR и оценке устойчивости к
дрейфу. По сравнению с лидером по F1 среди классических подходов
(XGBoost) целевая модель обеспечивает в одиннадцать раз меньшее
число ложных срабатываний на типовом потоке корпоративного SOC за
счет существенно более низкого FPR на операционном пороге.

Структура отчета. Основная часть отчета состоит из четырех
разделов, соответствующих четырем этапам, заявленным в постановке
задачи. Этап~1 содержит теоретический анализ методов обнаружения
аномалий, сравнительный анализ нейросетевых архитектур и публично
доступных датасетов NIDS, обзор методов активного тестирования,
формализацию требований и гипотезы. Этап~2 --- проектирование
подсистемы предобработки, архитектуры детектора, схемы калибровки
порога, формирования алертов, архитектуры конечно-автоматного
агента и инжектора дрейфа, а также формирование протокола
экспериментов. Этап~3 --- программная реализация прототипа
системы обнаружения, обучение и оптимизация моделей, разработка
конечно-автоматного агента. Этап~4 --- комплексное тестирование
прототипа, анализ эффективности, cross-evaluation, оценка
эксплуатационных характеристик и формирование итоговых материалов;
именно этому этапу уделено наибольшее внимание (расчет метрик,
анализ ошибок, тесты с конечно-автоматным агентом, оценка
устойчивости к дрейфу и производительности).
